\chapter{Unit 4: Number Systems}
\label{chap:unit4_qb}

\section{Practice Question Set (50 MCQs)}

\mcqitem{1}{What is the base (radix) of the Binary number system?}{2}{8}{10}{16}

\mcqitem{2}{What is the base (radix) of the Octal number system?}{2}{8}{10}{16}

\mcqitem{3}{What is the base (radix) of the Hexadecimal number system?}{8}{10}{12}{16}

\mcqitem{4}{Which digits are valid in the Binary number system?}{0 and 1 only}{0 through 7}{0 through 9}{0 through 9 and A through F}

\mcqitem{5}{Which of the following is an INVALID octal number?}{107}{543}{782}{615}

\mcqitem{6}{In the hexadecimal number system, which letter represents the decimal value 10?}{A}{B}{C}{D}

\mcqitem{7}{In the hexadecimal number system, which letter represents the decimal value 15?}{C}{D}{E}{F}

\mcqitem{8}{In the hexadecimal number system, what decimal value is represented by the letter 'B'?}{11}{12}{13}{14}

\mcqitem{9}{In the hexadecimal number system, what decimal value is represented by the letter 'D'?}{11}{12}{13}{14}

\mcqitem{10}{In positional number systems, what is the weight of the units position ($i=0$) in any base $b$?}{$0$}{$1$ ($b^0$)}{$b$}{$10$}

\mcqitem{11}{What is the decimal equivalent of the binary number $(1010)_2$?}{8}{10}{12}{14}

\mcqitem{12}{What is the decimal equivalent of the binary number $(1111)_2$?}{14}{15}{16}{17}

\mcqitem{13}{What is the decimal equivalent of the binary number $(1001)_2$?}{7}{8}{9}{10}

\mcqitem{14}{What is the binary representation of the decimal integer $13_{10}$?}{$(1101)_2$}{$(1011)_2$}{$(1110)_2$}{$(1001)_2$}

\mcqitem{15}{What is the binary representation of the decimal integer $8_{10}$?}{$(1000)_2$}{$(0100)_2$}{$(1100)_2$}{$(1111)_2$}

\mcqitem{16}{How many binary bits correspond directly to a single Octal digit?}{2 bits}{3 bits}{4 bits}{8 bits}

\mcqitem{17}{How many binary bits correspond directly to a single Hexadecimal digit?}{2 bits}{3 bits}{4 bits}{8 bits}

\mcqitem{18}{What is the octal equivalent of the binary number $(110101)_2$?}{$(55)_8$}{$(65)_8$}{$(35)_8$}{$(71)_8$}

\mcqitem{19}{What is the octal equivalent of the binary number $(101111)_2$?}{$(57)_8$}{$(47)_8$}{$(67)_8$}{$(55)_8$}

\mcqitem{20}{What is the binary equivalent of the octal number $(73)_8$?}{$(111011)_2$}{$(110011)_2$}{$(111010)_2$}{$(101011)_2$}

\mcqitem{21}{What is the hexadecimal equivalent of the binary number $(11001010)_2$?}{$(\text{CA})_{16}$}{$(\text{AC})_{16}$}{$(\text{BA})_{16}$}{$(\text{C}9)_{16}$}

\mcqitem{22}{What is the hexadecimal equivalent of the binary number $(10101111)_2$?}{$(\text{AF})_{16}$}{$(\text{FA})_{16}$}{$(\text{BF})_{16}$}{$(\text{AE})_{16}$}

\mcqitem{23}{What is the binary equivalent of the hexadecimal number $(\text{F}5)_{16}$?}{$(11110101)_2$}{$(11100101)_2$}{$(11110110)_2$}{$(10100101)_2$}

\mcqitem{24}{In binary addition, what is the sum of $1 + 1$?}{$0$ with no carry}{$0$ with a carry of $1$ (written as $10_2$)}{$1$ with no carry}{$11_2$}

\mcqitem{25}{In binary addition, what is the result of $1 + 1 + 1$?}{$10_2$}{$11_2$}{$01_2$}{$100_2$}

\mcqitem{26}{What is the decimal value of the binary fraction $(0.1)_2$?}{0.1}{0.5}{0.25}{0.125}

\mcqitem{27}{What is the decimal value of the binary fraction $(0.01)_2$?}{0.5}{0.25}{0.125}{0.01}

\mcqitem{28}{What is the decimal value of the binary number $(11.1)_2$?}{3.5}{3.25}{3.1}{2.5}

\mcqitem{29}{In a positional number, what is the leftmost non-zero digit called?}{Least Significant Digit (LSD)}{Most Significant Digit (MSD)}{Radix digit}{Fractional bit}

\mcqitem{30}{In a positional number, what is the rightmost digit called?}{Most Significant Digit (MSD)}{Least Significant Digit (LSD)}{Base exponent}{Binary root}

\mcqitem{31}{What is the separator between the integer part and the fractional part of a number in any base called?}{Comma}{Radix point}{Decimal mark}{Modulo}

\mcqitem{32}{When converting a decimal integer to another base by repeated division, how are the remainders recorded to form the answer?}{From first remainder to last remainder}{From last remainder (MSD) to first remainder (LSD)}{In numerical ascending order}{In numerical descending order}

\mcqitem{33}{When converting a decimal fraction to binary by repeated multiplication by 2, how are the generated integer parts recorded?}{From top (first generated) to bottom (last generated)}{From bottom to top}{In reverse order}{Alternating order}

\mcqitem{34}{What is the decimal equivalent of the octal number $(25)_8$?}{17}{21}{25}{29}

\mcqitem{35}{What is the decimal equivalent of the octal number $(70)_8$?}{56}{64}{70}{49}

\mcqitem{36}{What is the decimal equivalent of the hexadecimal number $(1\text{A})_{16}$?}{24}{26}{28}{30}

\mcqitem{37}{What is the decimal equivalent of the hexadecimal number $(20)_{16}$?}{20}{32}{40}{16}

\mcqitem{38}{Why is hexadecimal widely used by programmers to display memory addresses and machine code?}{Because computers use 16-state transistors}{Because one hex digit compactly represents exactly 4 binary bits}{Because hex eliminates all errors}{Because hex cannot store numbers}

\mcqitem{39}{In web design, what does the hexadecimal color code \texttt{\#FF0000} represent?}{Pure Blue}{Pure Green}{Pure Red}{Pure Yellow}

\mcqitem{40}{Unix-style file permissions such as \texttt{755} use octal notation where the permission value \texttt{7} represents which 3-bit binary pattern?}{$100_2$}{$110_2$}{$111_2$}{$011_2$}

\mcqitem{41}{Which of the following numbers is an INVALID hexadecimal numeral?}{2AF}{8B9}{19G}{C4E}

\mcqitem{42}{Convert decimal $156_{10}$ into its hexadecimal equivalent:}{$(\text{A}4)_{16}$}{$(9\text{C})_{16}$}{$(8\text{E})_{16}$}{$(9\text{D})_{16}$}

\mcqitem{43}{Convert decimal $156_{10}$ into its octal equivalent:}{$(234)_8$}{$(244)_8$}{$(224)_8$}{$(254)_8$}

\mcqitem{44}{What is the hexadecimal equivalent of the octal number $(75)_8$?}{First convert to binary $111101_2$, then group into $0011\;1101_2$ to get:}{$(3\text{D})_{16}$}{$(4\text{E})_{16}$}{$(2\text{C})_{16}$}{$(\text{F}5)_{16}$}

\mcqitem{45}{Why does a computer use binary internally rather than decimal?}{Binary circuits with two voltage states (ON/OFF) are simple, highly reliable, and noise-tolerant}{Binary was invented first by the Greeks}{Decimal numbers take up more space on disks}{Binary is easier for humans to read}

\mcqitem{46}{Why does the decimal fraction $0.1_{10}$ cause small rounding errors in standard binary computer arithmetic?}{Because $0.1$ is an irrational number}{Because $0.1_{10}$ results in an infinitely repeating binary fraction ($0.000110011\dots_2$)}{Because the CPU ignores fractions}{Because decimal is base 8}

\mcqitem{47}{Leading zeros added to the left of an integer (such as writing $00101_2$ instead of $101_2$):}{Double the value}{Do not change the mathematical value}{Halve the value}{Turn it negative}

\mcqitem{48}{A byte contains 8 bits. What is the maximum decimal value that can be represented by an unsigned 8-bit binary number ($11111111_2$)?}{128}{255}{256}{512}

\mcqitem{49}{What is the value of $(100)_2 + (011)_2$?}{$(111)_2$}{$(110)_2$}{$(101)_2$}{$(1000)_2$}

\mcqitem{50}{How is the hexadecimal prefix commonly written in programming languages like C, C++, and Python?}{\texttt{bin}}{\texttt{0x}}{\texttt{\#h}}{\texttt{dec}}

\newpage
\section{Answer Key \& Explanations --- Unit 4}

\begin{answerkeytable}{Unit 4: Number Systems --- Answer Key}
\begin{tabular}{*{10}{c}}
\toprule
\textbf{Q} & \textbf{Ans} & \textbf{Q} & \textbf{Ans} & \textbf{Q} & \textbf{Ans} & \textbf{Q} & \textbf{Ans} & \textbf{Q} & \textbf{Ans} \\
\midrule
1 & \textbf{A} & 11 & \textbf{B} & 21 & \textbf{A} & 31 & \textbf{B} & 41 & \textbf{C} \\
2 & \textbf{B} & 12 & \textbf{B} & 22 & \textbf{A} & 32 & \textbf{B} & 42 & \textbf{B} \\
3 & \textbf{D} & 13 & \textbf{C} & 23 & \textbf{A} & 33 & \textbf{A} & 43 & \textbf{A} \\
4 & \textbf{A} & 14 & \textbf{A} & 24 & \textbf{B} & 34 & \textbf{B} & 44 & \textbf{A} \\
5 & \textbf{C} & 15 & \textbf{A} & 25 & \textbf{B} & 35 & \textbf{A} & 45 & \textbf{A} \\
6 & \textbf{A} & 16 & \textbf{B} & 26 & \textbf{B} & 36 & \textbf{B} & 46 & \textbf{B} \\
7 & \textbf{D} & 17 & \textbf{C} & 27 & \textbf{B} & 37 & \textbf{B} & 47 & \textbf{B} \\
8 & \textbf{A} & 18 & \textbf{B} & 28 & \textbf{A} & 38 & \textbf{B} & 48 & \textbf{B} \\
9 & \textbf{C} & 19 & \textbf{A} & 29 & \textbf{B} & 39 & \textbf{C} & 49 & \textbf{A} \\
10 & \textbf{B} & 20 & \textbf{A} & 30 & \textbf{B} & 40 & \textbf{C} & 50 & \textbf{B} \\
\bottomrule
\end{tabular}
\end{answerkeytable}

\begin{expbox}[Selected Question Explanations]
\begin{itemize}[leftmargin=5mm]
    \item \textbf{Q.5 (C):} Octal digits range strictly from $0$ to $7$. The numeral $782$ contains '$8$', which is illegal in base 8.
    \item \textbf{Q.11 (B):} $1010_2 = 1(2^3) + 0(2^2) + 1(2^1) + 0(2^0) = 8 + 0 + 2 + 0 = 10_{10}$.
    \item \textbf{Q.14 (A):} $13 / 2 = 6$ rem $1$; $6/2 = 3$ rem $0$; $3/2 = 1$ rem $1$; $1/2 = 0$ rem $1$. Reading bottom to top gives $1101_2$.
    \item \textbf{Q.18 (B):} Grouping $110101_2$ in threes: $\underbrace{110}_{6}\;\underbrace{101}_{5} = (65)_8$.
    \item \textbf{Q.21 (A):} Grouping $11001010_2$ in fours: $\underbrace{1100}_{12=\text{C}}\;\underbrace{1010}_{10=\text{A}} = (\text{CA})_{16}$.
    \item \textbf{Q.26 (B):} $0.1_2 = 1 \times 2^{-1} = 1/2 = 0.5_{10}$.
    \item \textbf{Q.36 (B):} $(1\text{A})_{16} = 1(16^1) + 10(16^0) = 16 + 10 = 26_{10}$.
    \item \textbf{Q.42 (B):} $156 / 16 = 9$ with remainder $12$ ($\text{C}$). Thus, $156_{10} = (9\text{C})_{16}$.
    \item \textbf{Q.48 (B):} An 8-bit unsigned integer ranges from $0$ ($00000000_2$) to $2^8 - 1 = 255$ ($11111111_2$).
\end{itemize}
\end{expbox}

\newpage
